\section{Major Examination, 2016--17 (Semester 1)}

\noindent\textbf{Operating System- MTL458}\\
\textbf{Date:} 24-11-16,\quad
\textbf{Time:} 2 hour,\quad
\textbf{Total Marks:} 40

\begin{itemize}
  \item Write your Roll No.\ on the top immediately on receipt of the
        question paper.
  \item Attempt ALL questions along with any two parts from each
        question.  All questions carry equal marks.
\end{itemize}

\subsection*{Question 1}

\subsubsection*{(a) Address translation using the given page table}

\paragraph{Restatement.}
Suppose the page table for the process currently executing on the
processor looks like the following.  All numbers are decimal,
everything is numbered starting from zero, and all addresses are
memory byte addresses.  The page size is \(1024\) bytes.
\[
\begin{array}{c|c|c|c|c}
\mathrm{Virtual\ page\ number}&\mathrm{Valid\ bit}&\mathrm{Reference\ bit}&
\mathrm{Modify\ bit}&\mathrm{Page\ frame\ number}\\ \hline
0&1&1&0&4\\
1&1&1&1&7\\
2&0&0&0&-\\
3&1&0&0&2\\
4&0&0&0&-\\
5&1&0&1&0
\end{array}
\]

\paragraph{(a.1) Restatement.}
Describe exactly how, in general, a virtual address generated by the
CPU is translated into a physical main memory address.

\paragraph{Solution.}
For a page size \(P=1024\), the memory-management unit divides a virtual
address \(v\) into
\[
 \mathrm{VPN}=\left\lfloor\frac{v}{1024}\right\rfloor,
 \qquad
 d=v\bmod 1024,
\]
where \(d\) is the ten-bit page offset.  The VPN is looked up in the
TLB.  On a TLB miss, the hardware (or an operating-system handler,
depending on the architecture) uses the page-table base register and
the VPN to locate the process's page-table entry (PTE).

The PTE is then checked.  If the entry is outside the page table, the
virtual address is illegal.  If it exists but its valid/present bit is
zero, the access causes a page fault (or an illegal-address fault if the
page is not part of the process's address space).  Access-protection
bits are also checked.  For a valid and permitted entry containing
physical frame number \(f\), translation preserves the offset:
\[
 \mathrm{PA}=f(1024)+d.
\]
The reference bit is set by an access, and the modify/dirty bit is set
by a write.

\paragraph{(a.2) Restatement.}
What physical address, if any, would each of the following virtual
addresses correspond to?  (i) \(1052\), (ii) \(2221\), (iii) \(5499\).

\paragraph{Solution.}
The divisions by \(1024\) are
\[
\begin{array}{c|c|c|c}
v & \left\lfloor v/1024\right\rfloor & v\bmod 1024 & \mathrm{result}\\ \hline
1052 & 1 & 28 & 7(1024)+28=7196\\
2221 & 2 & 173 & \mathrm{page\ fault\ (valid\ bit}=0)\\
5499 & 5 & 379 & 0(1024)+379=379
\end{array}
\]
Thus \(1052\) maps through valid VPN \(1\), whose page frame number is
\(7\), and \(5499\) maps through valid VPN \(5\), whose page frame
number is \(0\).  VPN \(2\), selected by \(2221\), has its valid bit
clear, so that reference produces a page fault and has no physical
address until the fault is serviced.

\subsubsection*{(b) Number of page-table levels}

\paragraph{Restatement.}
Assuming a page size of \(4\) Kbytes and that a page table entry takes
\(4\) bytes, how many levels of page tables would be required to map a
\(64\)-bit address space, if the top level page table fits into a single
page?

\paragraph{Solution.}
A \(4\,\mathrm{KB}=2^{12}\)-byte page gives a \(12\)-bit offset, leaving
\[
 64-12=52
\]
virtual-page-number bits.  One page-table page holds
\[
 \frac{2^{12}}{2^2}=2^{10}=1024
\]
PTEs, so a full page-table level can consume at most \(10\) VPN bits.
Thus
\[
 L=\left\lceil\frac{52}{10}\right\rceil=6.
\]
The address split, from most significant to least significant bits, is
\[
 \underbrace{2}_{L_1}\;|\;
 \underbrace{10}_{L_2}\;|\;
 \underbrace{10}_{L_3}\;|\;
 \underbrace{10}_{L_4}\;|\;
 \underbrace{10}_{L_5}\;|\;
 \underbrace{10}_{L_6}\;|\;
 \underbrace{12}_{\mathrm{offset}}.
\]
Hence six levels are required.  The top level uses only a two-bit
index (four of the \(1024\) possible PTE positions), and therefore
comfortably fits in one page.

\subsubsection*{(c) LRU page replacement}

\paragraph{Restatement.}
Consider the following page reference strings:
\[
1,2,3,4,2,1,5,6,2,1,2,3,7,6,3,2,1,2,3,6.
\]
How many page faults would occur for the LRU replacement algorithm,
assuming four frames?

\paragraph{Solution.}
The four frame columns below denote fixed frame slots, not recency
order.  A dash in the victim column means either a hit or use of a
previously empty frame.
\[
\begin{array}{c|c|c|c|c|c|c|c}
\mathrm{Step}&\mathrm{Ref.}&F_1&F_2&F_3&F_4&\mathrm{Result}&\mathrm{Victim}\\ \hline
1 &1&1&-&-&-&\mathrm{fault}&-\\
2 &2&1&2&-&-&\mathrm{fault}&-\\
3 &3&1&2&3&-&\mathrm{fault}&-\\
4 &4&1&2&3&4&\mathrm{fault}&-\\
5 &2&1&2&3&4&\mathrm{hit}&-\\
6 &1&1&2&3&4&\mathrm{hit}&-\\
7 &5&1&2&5&4&\mathrm{fault}&3\\
8 &6&1&2&5&6&\mathrm{fault}&4\\
9 &2&1&2&5&6&\mathrm{hit}&-\\
10&1&1&2&5&6&\mathrm{hit}&-\\
11&2&1&2&5&6&\mathrm{hit}&-\\
12&3&1&2&3&6&\mathrm{fault}&5\\
13&7&1&2&3&7&\mathrm{fault}&6\\
14&6&6&2&3&7&\mathrm{fault}&1\\
15&3&6&2&3&7&\mathrm{hit}&-\\
16&2&6&2&3&7&\mathrm{hit}&-\\
17&1&6&2&3&1&\mathrm{fault}&7\\
18&2&6&2&3&1&\mathrm{hit}&-\\
19&3&6&2&3&1&\mathrm{hit}&-\\
20&6&6&2&3&1&\mathrm{hit}&-
\end{array}
\]
Faults occur at steps
\[
1,2,3,4,7,8,12,13,14,17,
\]
so the exact number of page faults is
\[
\boxed{10}.
\]

\subsection*{Question 2}

\subsubsection*{(a) Banker's safety test}

\paragraph{Restatement.}
Consider a system with a total of \(150\) units of memory, allocated to
three processes as shown:
\[
\begin{array}{c|c|c}
\mathrm{Process}&\mathrm{Max}&\mathrm{Hold}\\ \hline
1&70&45\\
2&60&40\\
3&60&15
\end{array}
\]
Apply the banker's algorithm to determine whether it would be safe to
grant each of the following requests.  If yes, indicate a sequence of
terminations that could be guaranteed possible.  If no, show the
reduction of the resulting allocation table.

\begin{itemize}
  \item[(a.1)] A fourth process arrives, with a maximum memory need of
        \(60\) and an initial need of \(25\) units.
  \item[(a.2)] A fourth process arrives, with a maximum memory need of
        \(60\) and an initial need of \(35\) units.
\end{itemize}

\paragraph{Initial state.}
The first three processes hold
\[
45+40+15=100
\]
units, so before admitting \(P_4\) the available amount is
\[
150-100=50.
\]

\paragraph{(a.1) Grant \(25\) units to \(P_4\).}
After the grant, \(P_4\)'s remaining need is \(60-25=35\), and
\[
\mathrm{Work}=50-25=25.
\]
A valid safety reduction is
\[
\begin{array}{c|c|c|c}
\mathrm{Chosen\ process}&\mathrm{Need}&
\mathrm{Allocation\ released}&\mathrm{Work\ after\ completion}\\ \hline
P_1&25&45&25+45=70\\
P_2&20&40&70+40=110\\
P_3&45&15&110+15=125\\
P_4&35&25&125+25=150
\end{array}
\]
At every row, the process's remaining need is no greater than the
current Work.  Therefore the grant is safe; one safe sequence is
\[
\boxed{\langle P_1,P_2,P_3,P_4\rangle}.
\]

\paragraph{(a.2) Grant \(35\) units to \(P_4\).}
Now \(P_4\)'s remaining need is \(60-35=25\), while
\[
\mathrm{Work}=50-35=15.
\]
The first reduction step already fails:
\[
\begin{array}{c|c|c|c}
\mathrm{Process}&\mathrm{Allocation}&\mathrm{Need}&
\mathrm{Need}\leq 15?\\ \hline
P_1&45&25&\mathrm{No}\\
P_2&40&20&\mathrm{No}\\
P_3&15&45&\mathrm{No}\\
P_4&35&25&\mathrm{No}
\end{array}
\]
No process can acquire its remaining claim and finish, so Work remains
\(15\) and all four processes remain unfinished.  The resulting state
is unsafe (and, if all processes make their full remaining requests,
deadlocked).  The request must not be granted.

\subsubsection*{(b) Sufficient condition excluding deadlock}

\paragraph{Restatement.}
\(N\) processes share \(M\) resource units that can be reserved and
released only one at a time.  The maximum need of each process does not
exceed \(M\), and the sum of all maximum needs is less than \(M+N\).
Show that a deadlock can not occur.

\paragraph{Solution.}
Let \(a_i\) be the current allocation of process \(i\), let
\(r_i=m_i-a_i\) be its maximum remaining need, and let \(A\) be the
number of available units.  For \(n\) currently unfinished processes,
\[
\sum_i a_i=M-A.
\]
Suppose, for contradiction, that none of these processes can have its
remaining need met from the available resources.  Since resource
counts are integers, this says \(r_i>A\), or \(r_i\geq A+1\), for every
unfinished process.  It follows that
\[
\begin{split}
\sum_i m_i
 &=\sum_i a_i+\sum_i r_i\\
 &\geq (M-A)+n(A+1)\\
 &=M+n+(n-1)A\\
 &\geq M+n,
\end{split}
\]
contradicting \(\sum_i m_i<M+n\).  Thus at least one process has
\(r_i\leq A\); its entire remaining claim can be satisfied, after which
it terminates and releases its allocation.

The condition is preserved when a process terminates.  Because all
claims are positive integers, from
\(\sum_i m_i\leq M+n-1\), removing a process of claim \(m_j\geq1\)
leaves
\[
\sum_{i\ne j}m_i\leq M+n-2<M+(n-1).
\]
The argument can therefore be repeated until every process finishes.
Equivalently, at a putative deadlock every unfinished process would
hold at most one less than its maximum, but the stated inequality
ensures that the system cannot exhaust all \(M\) units in that
configuration.  Hence deadlock is impossible (assuming requests that
can be granted are not postponed forever).

\subsubsection*{(c) Lefty/righty dining philosophers}

\paragraph{Restatement.}
Suppose that there are two types of philosophers.  One type always
picks up his left fork first (a ``lefty''), and the other type always
picks up his right fork first (a ``righty'').  Using semaphore, define
both types of philosophers' behavior.

\paragraph{Solution.}
For \(n\) philosophers, let each fork be a binary semaphore:
\[
\mathrm{semaphore}\;\mathrm{fork}[0\ldots n-1],
\qquad \mathrm{fork}[k]=1\quad\hbox{initially}.
\]
For philosopher \(i\), define
\(\mathrm{left}(i)=i\) and
\(\mathrm{right}(i)=(i+1)\bmod n\).  Here \(P\) is wait/down and \(V\)
is signal/up.

\begin{verbatim}
lefty(i):
    while (true) {
        think();
        P(fork[left(i)]);
        P(fork[right(i)]);
        eat();
        V(fork[right(i)]);
        V(fork[left(i)]);
    }

righty(i):
    while (true) {
        think();
        P(fork[right(i)]);
        P(fork[left(i)]);
        eat();
        V(fork[left(i)]);
        V(fork[right(i)]);
    }
\end{verbatim}

Each fork is used by at most one philosopher because its semaphore is
binary.  If the circle contains both acquisition orientations (the
usual implementation makes one philosopher a righty and all the
others lefties), the uniform circular-wait pattern is broken, so all
philosophers cannot deadlock while holding their first fork.  If every
philosopher has the same orientation, the familiar circular deadlock
is still possible.

\subsection*{Question 3}

\subsubsection*{(a) All concurrent traces}

\paragraph{Restatement.}
Consider a concurrent program with two processes, \(p\) and \(q\),
defined as follows.  \(A,B,C,D,\) and \(E\) are arbitrary atomic
(indivisible) statements.  Assume that the main program (not shown)
does a \texttt{parbegin} of the two processes.
\begin{verbatim}
Void p()             void q()
{ A; B; C; }         { D; E; }
\end{verbatim}
Show all the possible interleavings of the execution of the preceding
two processes (show this by giving execution ``traces'' in terms of the
atomic statements).

\paragraph{Solution.}
Every trace must preserve \(A<B<C\) and \(D<E\).  Choosing the three
positions occupied by \(A,B,C\) gives
\[
\binom{5}{3}=\binom{5}{2}=10
\]
traces.  All ten are
\[
\begin{array}{lllll}
ABCDE, & ABDCE, & ABDEC, & ADBCE, & ADBEC,\\
ADEBC, & DABCE, & DABEC, & DAEBC, & DEABC.
\end{array}
\]

\subsubsection*{(b) Peterson's algorithm}

\paragraph{Restatement.}
Prove the correctness of Peterson's algorithm for mutual exclusion and
progress of some processes requiring access to its critical section.

\paragraph{Solution.}
The shared variables are initialized by
\[
\mathrm{flag}[0]=\mathrm{flag}[1]=\mathrm{false};
\]
\(\mathrm{turn}\) may initially have either value.  Process \(P_i\),
where \(j=1-i\), executes:

\begin{verbatim}
while (true) {
    flag[i] = true;
    turn = j;
    while (flag[j] && turn == j)
        ;                       /* spin */

    critical_section();

    flag[i] = false;
    remainder_section();
}
\end{verbatim}

This proof assumes atomic reads and writes, sequentially consistent
shared memory, and that a process does not remain in its critical
section forever.

\paragraph{Mutual exclusion.}
Assume both \(P_0\) and \(P_1\) are simultaneously in their critical
sections.  Both flags must then be true.  Consider the two entry-section
writes to \(\mathrm{turn}\), and suppose without loss of generality that
\(P_0\)'s write \(\mathrm{turn}=1\) occurs last.  Before that write,
\(P_1\) has already set \(\mathrm{flag}[1]\) to true, and that flag
cannot be cleared before \(P_1\) leaves its critical section.  Thus
\(P_0\) subsequently tests
\(\mathrm{flag}[1]\mathbin{\&\&}\mathrm{turn}=1\), which is true, and
cannot enter while \(P_1\) is in its critical section.  This
contradicts the assumption that both entered.  The case in which
\(P_1\)'s turn write is last is symmetric.  Therefore at most one
process is in the critical section.

\paragraph{Progress and bounded waiting.}
If only \(P_i\) wants to enter, then \(\mathrm{flag}[j]\) is false and
\(P_i\) enters immediately.  If both want to enter, both set their
flags, and the last write to \(\mathrm{turn}\) determines which process
waits.  Exactly one loop condition can remain true, so the other
process enters; there is no state in which both wait for each other.
On leaving, that process clears its flag, enabling the waiting process.
Moreover, after \(P_i\) announces its interest, \(P_j\) can enter ahead
of it at most once before \(P_i\) enters.  Thus the algorithm supplies
mutual exclusion, progress, and bounded waiting under fair execution.

\subsubsection*{(c) An incorrect bounded-buffer implementation}

\paragraph{Restatement.}
Provide the algorithms for producer and consumer processes in bounded
buffer problem and show that they may not function correctly when
executed concurrently.

\paragraph{Solution.}
The following tempting implementation uses a circular buffer of size
\(N\) and a shared count, but no atomic operation, lock, or semaphore:

\begin{verbatim}
/* shared: item buffer[N]; int in = 0, out = 0, count = 0; */

produce(x):
    while (count == N)
        ;                       /* wait for a free slot */
    buffer[in] = x;
    in = (in + 1) % N;
    count = count + 1;

consume():
    while (count == 0)
        ;                       /* wait for an item */
    x = buffer[out];
    out = (out + 1) % N;
    count = count - 1;
    return x;
\end{verbatim}

The source-level increments are read--modify--write sequences.  For
example, suppose initially \(\mathrm{count}=1\), and a producer and a
consumer run concurrently:
\[
\begin{array}{c|l}
1&\mathrm{producer\ reads\ count}=1\\
2&\mathrm{consumer\ reads\ count}=1\\
3&\mathrm{producer\ computes}\ 2\\
4&\mathrm{consumer\ computes}\ 0\\
5&\mathrm{producer\ stores}\ 2\\
6&\mathrm{consumer\ stores}\ 0
\end{array}
\]
One item was added and one removed, so the correct final count is
\(1\), but the stored count is \(0\).  Reversing the final two stores
makes it \(2\), which is also wrong.  Future executions can therefore
underflow, overflow, or wait despite the actual buffer state.

With multiple producers, two producers can also both observe
\(\mathrm{count}=N-1\), pass the full test, use the same value of
\(\mathrm{in}\), and overwrite one another's item.  Analogous races
exist for multiple consumers at \(\mathrm{out}\).  Correct bounded
buffer code needs mutual exclusion around buffer/index updates and
counting synchronization, conventionally semaphores
\(\mathrm{empty}=N\), \(\mathrm{full}=0\), and
\(\mathrm{mutex}=1\).

\subsection*{Question 4}

\subsubsection*{(a) Optimality of nonpreemptive SPN/SJF}

\paragraph{Restatement.}
Prove that, among nonpreemptive scheduling algorithm, SPN (shortest
process next) provides the minimum average waiting time for a batch of
jobs that arrive at the same time.  Assume that the scheduler must
always execute a task if one is available.

\paragraph{Solution.}
Consider any schedule containing two adjacent jobs in the order
\(J_x,J_y\), with service times \(x>y\).  Let \(T\) be the total service
time before these two jobs.  Their waiting times in the original order
are
\[
W_x=T,\qquad W_y=T+x,
\]
so their contribution to total waiting time is \(2T+x\).  After
swapping them, it is
\[
W'_y=T,\qquad W'_x=T+y,
\]
whose sum is \(2T+y\).  The swap reduces total waiting time by
\[
(2T+x)-(2T+y)=x-y>0.
\]
No earlier job is affected, and later jobs start after the same total
amount \(x+y\), so their waiting times do not change.

Thus any schedule containing a longer job immediately before a shorter
one can be strictly improved.  Repeatedly exchanging such inversions
produces nondecreasing service-time order, exactly SPN/SJF.  Dividing
the minimum total waiting time by the fixed number of jobs also
minimizes average waiting time.  The work-conserving condition matters:
with all jobs already available, deliberate idle time can only add
waiting and cannot improve this objective.

\subsubsection*{(b) SRTF schedule and turnaround times}

\paragraph{Restatement.}
Consider the set of any four processes, with the length of the CPU
burst \((8,4,9,5)\) and arrival time \((0,1,2,3)\) given in
milliseconds:
\[
\begin{array}{c|c|c}
\mathrm{Process}&\mathrm{Arrival\ (ms)}&\mathrm{CPU\ burst\ (ms)}\\ \hline
P_1&0&8\\
P_2&1&4\\
P_3&2&9\\
P_4&3&5
\end{array}
\]
Draw a Gantt chart that illustrates the execution of these processes
using shortest-remaining-time-first scheduling.  Also, calculate
turnaround time of each process.

\paragraph{Solution.}
At time \(0\), only \(P_1\) is ready.  At time \(1\), \(P_1\) has
seven milliseconds remaining and the newly arrived \(P_2\) needs only
four, so \(P_2\) preempts \(P_1\).  The arrivals at times \(2\) and
\(3\) do not preempt \(P_2\), whose remaining times are then three and
two milliseconds.  After \(P_2\) finishes, the shortest remaining jobs
are \(P_4\) (five), \(P_1\) (seven), and \(P_3\) (nine), in that order.
The Gantt chart is
\[
\begin{array}{c|ccccc}
\mathrm{time\ interval}&[0,1)&[1,5)&[5,10)&[10,17)&[17,26)\\ \hline
\mathrm{running}&P_1&P_2&P_4&P_1&P_3
\end{array}
\]
Thus
\[
C_1=17,\qquad C_2=5,\qquad C_3=26,\qquad C_4=10.
\]
Using \(T_i=C_i-A_i\),
\[
\begin{array}{c|c|c|c|c}
\mathrm{Process}&C_i&A_i&T_i=C_i-A_i&
W_i=T_i-\mathrm{burst}_i\\ \hline
P_1&17&0&17&9\\
P_2&5&1&4&0\\
P_3&26&2&24&15\\
P_4&10&3&7&2
\end{array}
\]
Therefore the turnaround times are
\[
\boxed{T_1=17,\quad T_2=4,\quad T_3=24,\quad T_4=7\ \mathrm{ms}},
\]
and their average is
\[
\frac{17+4+24+7}{4}=\frac{52}{4}=13\ \mathrm{ms}.
\]

\subsubsection*{(c) Criteria for comparing CPU schedulers}

\paragraph{Restatement.}
Suggest a criteria for comparing CPU-scheduling algorithms.

\paragraph{Solution.}
Important criteria include:
\begin{itemize}
  \item \emph{CPU utilization}: the fraction of elapsed time for which
        the CPU does useful work; maximize it.
  \item \emph{Throughput}: completed processes per unit time; maximize
        it.
  \item \emph{Turnaround time}: completion time minus arrival time;
        minimize its mean and, where relevant, its tail.
  \item \emph{Waiting time}: total time spent in ready queues; minimize
        it.
  \item \emph{Response time}: time from arrival/request until the first
        CPU service or visible response; especially important for
        interactive systems.
  \item \emph{Fairness and starvation freedom}: give processes a
        reasonable share while respecting priorities; bounded waiting
        is desirable.
  \item \emph{Predictability and deadline satisfaction}: low variance
        and meeting real-time deadlines may matter more than the mean.
  \item \emph{Scheduling overhead}: context switches, migrations, and
        scheduler computation should consume little CPU time; affinity
        and cache behavior also matter on multiprocessors.
\end{itemize}
The appropriate choice depends on workload: for example, a batch
system often emphasizes throughput and turnaround, whereas an
interactive system emphasizes response time and fairness.
